%%======================================================================
%%  Conclusion
%%======================================================================
\chapter*{Conclusion}
\addcontentsline{toc}{chapter}{Conclusion}

Au terme de ce stage, nous avons conçu et développé une \textbf{plateforme éducative web complète} répondant aux besoins de gestion d'un établissement d'enseignement supérieur. Le projet \textit{EduPlatform} couvre l'ensemble du cycle de vie pédagogique et administratif, de la gestion des cours et des devoirs au suivi des absences, en passant par le traitement des demandes et des réclamations.

\vspace{0.5cm}

\textbf{Bilan technique :}

\begin{itemize}[noitemsep]
    \item Un \textbf{backend Spring Boot} robuste composé de 25 entités JPA, 28 contrôleurs REST et plus de 90 endpoints, structuré selon le pattern MVC ;
    \item Un \textbf{frontend React/TypeScript} réactif comprenant plus de 30 pages spécialisées, avec trois espaces dédiés (Administrateur, Enseignant, Étudiant) ;
    \item Une \textbf{base de données relationnelle} bien normalisée avec une stratégie d'héritage SINGLE\_TABLE et 23 tables ;
    \item Un système de \textbf{notifications intelligentes} avec workflow d'approbation et rappels automatiques ;
    \item Un \textbf{workflow de réclamation de notes} multi-étapes impliquant trois acteurs ;
    \item Un \textbf{système d'upload de fichiers} pour les ressources de cours, les soumissions et les justifications.
\end{itemize}

\vspace{0.5cm}

\textbf{Compétences acquises :}

Ce stage nous a permis d'acquérir et de consolider de nombreuses compétences :
\begin{itemize}[noitemsep]
    \item Maîtrise du framework \textbf{Spring Boot} et de l'écosystème Java pour le développement backend ;
    \item Développement d'interfaces modernes avec \textbf{React}, \textbf{TypeScript} et \textbf{Tailwind CSS} ;
    \item Conception de bases de données relationnelles et utilisation avancée de \textbf{JPA/Hibernate} ;
    \item Modélisation UML (diagrammes de classes, de cas d'utilisation, de séquence) ;
    \item Pratique de l'architecture REST et communication client-serveur ;
    \item Gestion de projet et organisation du code en couches.
\end{itemize}

\vspace{0.5cm}

\textbf{Perspectives d'amélioration :}

Plusieurs pistes d'évolution sont envisageables pour renforcer la plateforme :
\begin{itemize}[noitemsep]
    \item \textbf{Sécurisation de l'API} : mise en place de Spring Security avec authentification JWT côté serveur ;
    \item \textbf{Tests unitaires et d'intégration} : ajout de tests automatisés (JUnit, Mockito, React Testing Library) ;
    \item \textbf{Conteneurisation} : déploiement via Docker et Docker Compose pour faciliter la mise en production ;
    \item \textbf{Notifications en temps réel} : intégration de WebSocket (STOMP) pour les notifications push ;
    \item \textbf{Tableau de bord analytique avancé} : graphiques interactifs supplémentaires, export PDF des relevés de notes ;
    \item \textbf{Application mobile} : développement d'une application mobile (React Native) pour un accès en mobilité.
\end{itemize}

\vspace{0.5cm}

En définitive, ce projet de stage a été une expérience enrichissante qui nous a permis de mettre en pratique les connaissances théoriques acquises durant notre formation, tout en découvrant les réalités du développement logiciel dans un contexte professionnel.

\newpage
